% problem_id: S001-lemma-directed-colimit-diagram-finite-presentation
% source_id: S001 (Stacks Project)
% source_file: properties.tex
% statement_locator: properties.tex lines 3039--3064
% proof_locator: properties.tex lines 3066--3190
% env: lemma
% id_note: label 在本来源内唯一
% pairing: tight (verified)
% status: complete (statement + author proof verbatim + reason + locators)
%
\begin{lemma}
\label{lemma-directed-colimit-diagram-finite-presentation}
\begin{slogan}
Quasi-coherent modules on quasi-compact and quasi-separated schemes
are filtered colimits of finitely presented modules.
\end{slogan}
Let $X$ be a scheme. Assume $X$ is quasi-compact and quasi-separated.
Let $\mathcal{F}$ be a quasi-coherent $\mathcal{O}_X$-module.
There exist
\begin{enumerate}
\item a filtered index category $\mathcal{I}$ (see
Categories, Definition \ref{categories-definition-directed}),
\item a diagram $\mathcal{I} \to \textit{Mod}(\mathcal{O}_X)$ (see
Categories, Section \ref{categories-section-limits}),
$i \mapsto \mathcal{F}_i$,
\item morphisms of $\mathcal{O}_X$-modules
$\varphi_i : \mathcal{F}_i \to \mathcal{F}$
\end{enumerate}
such that each $\mathcal{F}_i$ is of finite presentation
and such that the morphisms $\varphi_i$ induce an isomorphism
$$
\colim_i \mathcal{F}_i
=
\mathcal{F}.
$$
\end{lemma}

\begin{proof}
Choose a set $I$ and for each $i \in I$ an $\mathcal{O}_X$-module
of finite presentation and a homomorphism of $\mathcal{O}_X$-modules
$\varphi_i : \mathcal{F}_i \to \mathcal{F}$ with the following
property: For any $\psi : \mathcal{G} \to \mathcal{F}$ with $\mathcal{G}$
of finite presentation there is an $i \in I$ such that there exists
an isomorphism $\alpha : \mathcal{F}_i \to \mathcal{G}$ with
$\varphi_i = \psi \circ \alpha$. It is clear from
Modules, Lemma \ref{modules-lemma-set-isomorphism-classes-finite-type-modules}
that such a set exists (see also its proof).
We denote $\mathcal{I}$ the category
with $\Ob(\mathcal{I}) = I$ and given $i, i' \in I$
we set
$$
\Mor_\mathcal{I}(i, i') =
\{\alpha : \mathcal{F}_i \to \mathcal{F}_{i'} \mid
\alpha \circ \varphi_{i'} = \varphi_i
\}.
$$
We claim that $\mathcal{I}$ is a filtered category and that
$\mathcal{F} = \colim_i \mathcal{F}_i$.

\medskip\noindent
Let $i, i' \in I$. Then we can consider the morphism
$$
\mathcal{F}_i \oplus \mathcal{F}_{i'} \longrightarrow \mathcal{F}
$$
which is the direct sum of $\varphi_i$ and $\varphi_{i'}$.
Since a direct sum of finitely presented $\mathcal{O}_X$-modules
is finitely presented we see that there exists some $i'' \in I$
such that $\varphi_{i''} : \mathcal{F}_{i''} \to \mathcal{F}$
is isomorphic to the displayed arrow towards $\mathcal{F}$ above.
Since there are commutative diagrams
$$
\xymatrix{
\mathcal{F}_i \ar[r] \ar[d] & \mathcal{F} \ar@{=}[d] \\
\mathcal{F}_i \oplus \mathcal{F}_{i'} \ar[r] & \mathcal{F}
}
\quad
\text{and}
\quad
\xymatrix{
\mathcal{F}_{i'} \ar[r] \ar[d] & \mathcal{F} \ar@{=}[d] \\
\mathcal{F}_i \oplus \mathcal{F}_{i'} \ar[r] & \mathcal{F}
}
$$
we see that there are morphisms $i \to i''$ and $i' \to i''$
in $\mathcal{I}$. Next, suppose that we have $i, i' \in I$ and
morphisms $\alpha, \beta : i \to i'$ (corresponding to $\mathcal{O}_X$-module
maps $\alpha, \beta : \mathcal{F}_i \to \mathcal{F}_{i'}$).
In this case consider the coequalizer
$$
\mathcal{G} =
\Coker(
\mathcal{F}_i \xrightarrow{\alpha - \beta} \mathcal{F}_{i'}
)
$$
Note that $\mathcal{G}$ is an $\mathcal{O}_X$-module of finite presentation.
Since by definition of morphisms in the category $\mathcal{I}$
we have $\varphi_{i'} \circ \alpha = \varphi_{i'} \circ \beta$
we see that we get an induced map $\psi : \mathcal{G} \to \mathcal{F}$.
Hence again the pair $(\mathcal{G}, \psi)$ is isomorphic to
the pair $(\mathcal{F}_{i''}, \varphi_{i''})$ for some $i''$.
Hence we see that there exists a morphism $i' \to i''$ in
$\mathcal{I}$ which equalizes $\alpha$ and $\beta$. Thus we have
shown that the category $\mathcal{I}$ is filtered.

\medskip\noindent
We still have to show that the colimit of the diagram is $\mathcal{F}$.
By definition of the colimit, and by our definition of the category
$\mathcal{I}$ there is a canonical map
$$
\varphi :
\colim_i \mathcal{F}_i
\longrightarrow
\mathcal{F}.
$$
Pick $x \in X$. Let us show that $\varphi_x$ is an isomorphism.
Recall that
$$
(\colim_i \mathcal{F}_i)_x
=
\colim_i \mathcal{F}_{i, x},
$$
see
Sheaves, Section \ref{sheaves-section-limits-sheaves}.
First we show that the map $\varphi_x$ is injective.
Suppose that $s \in \mathcal{F}_{i, x}$ is an element
such that $s$ maps to zero in $\mathcal{F}_x$. Then there exists
a quasi-compact open $U$ such that $s$ comes from $s \in \mathcal{F}_i(U)$
and such that $\varphi_i(s) = 0$ in $\mathcal{F}(U)$.
By Lemma \ref{lemma-extend}
we can find a finite type quasi-coherent subsheaf
$\mathcal{K} \subset \Ker(\varphi_i)$ which restricts to
the quasi-coherent $\mathcal{O}_U$-submodule of $\mathcal{F}_i$
generated by $s$:
$\mathcal{K}|_U = \mathcal{O}_U\cdot s \subset \mathcal{F}_i|_U$.
Clearly, $\mathcal{F}_i/\mathcal{K}$ is of finite presentation and
the map $\varphi_i$ factors through the quotient map
$\mathcal{F}_i \to \mathcal{F}_i/\mathcal{K}$. Hence we can find
an $i' \in I$ and a morphism $\alpha : \mathcal{F}_i \to \mathcal{F}_{i'}$
in $\mathcal{I}$ which can be identified with the quotient map
$\mathcal{F}_i \to \mathcal{F}_i/\mathcal{K}$. Then it follows
that the section $s$ maps to zero in $\mathcal{F}_{i'}(U)$ and
in particular in
$(\colim_i \mathcal{F}_i)_x =
\colim_i \mathcal{F}_{i, x}$.
The injectivity follows.
Finally, we show that the map $\varphi_x$ is surjective.
Pick $s \in \mathcal{F}_x$. Choose a quasi-compact open neighbourhood
$U \subset X$ of $x$ such that $s$ corresponds to a section
$s \in \mathcal{F}(U)$. Consider the map
$s : \mathcal{O}_U \to \mathcal{F}$ (multiplication by $s$).
By Lemma \ref{lemma-extend-finite-presentation}
there exists an $\mathcal{O}_X$-module $\mathcal{G}$
of finite presentation and an $\mathcal{O}_X$-module map
$\mathcal{G} \to \mathcal{F}$ such that $\mathcal{G}|_U \to \mathcal{F}|_U$
is identified with
$s : \mathcal{O}_U \to \mathcal{F}$.
Again by definition of $\mathcal{I}$ there exists an $i \in I$
such that $\mathcal{G} \to \mathcal{F}$ is isomorphic to
$\varphi_i : \mathcal{F}_i \to \mathcal{F}$. Clearly there exists
a section $s' \in \mathcal{F}_i(U)$ mapping to $s \in \mathcal{F}(U)$.
This proves surjectivity and the proof of the lemma is complete.
\end{proof}
